\begin{textsample}{NEG Dim 3 – 2023\_c6 – Score 1.00 – 2023\_c6/enade\_2023\_eng\_mec\_63b17ce9db.txt}  \label{ex:f3_neg_184}
CAPÍTULO 5 ANÁLISE TÉCNICA DA PROVA Este capítulo tem o objetivo de apresentar o desempenho dos estudantes concluintes de Engenharia Mecânica no Enade 2023.
Para isso, foram calculadas as estatísticas básicas da prova em seu todo ( subseção 5.1.1 ), bem como as estatísticas relacionadas aos componentes de Formação Geral ( subseção 5.1.2 ) e de Conhecimento Específico da área de Engenharia Mecânica ( subseção 5.1.3 ).
Dadas as suas características, foram analisadas, em separado, as questões objetivas ( seção 5.2 ) e as questões discursivas ( seção 5.3 ).
Tomando -se como base a questão discursiva do componente de Formação Geral, nas subseções de 5.3.1.4 a 5.3.1.6, são apresentados os resultados e os comentários sobre a correção das respostas em relação ao desempenho linguístico.
Para as questões objetivas, foram disponibilizados os índices de facilidade e de discriminação ponto-bisserial, também separadamente, para os componentes de Formação Geral ( subseção 5.2.1 ) e de Conhecimento Específico ( subseção 5.2.2 ).
De cada componente, uma das questões foi escolhida para exemplificar a análise gráfica, relacionando -se as alternativas escolhidas pelos estudantes ( inclusive o gabarito ) ao número de acertos no componente.
No Anexo IV, apresenta -se a íntegra da análise gráfica para todas as questões objetivas.
Para cada uma das questões discursivas, os conteúdos dos tipos mais comuns de respostas dos estudantes são apresentados e comparados com o padrão de resposta esperado ( ver Anexo VIII ).
Nas tabelas, constam as seguintes estatísticas das notas5 : média do desempenho na prova, erro-padrão da média, desvio-padrão, nota mínima, mediana e nota máxima para cada um de seus componentes.
Tais estatísticas contemplam o total de estudantes concluintes da área de Engenharia Mecânica inscritos e presentes na prova do Enade 2023, tendo em vista agregações, ou por grandes regiões e o país como um todo, ou por categoria administrativa, organização acadêmica da IES e modalidade de oferta dos cursos.
Em relação aos gráficos de distribuição de notas, o intervalo considerado foi de 10 unidades, aberto à esquerda e fechado à direita, com exceção do primeiro intervalo, [ 0 ; 10 ], fechado em ambos os extremos.
Para os gráficos de distribuição das notas das questões discursivas, foram consideradas mais duas categorias : questão em branco6 e nota zero. 5 Uma definição dessas estatísticas pode ser encontrada no Glossário inserido no final desse relatório. 6 Nesse grupo estão incluídas também as respostas classificadas como nulas ou desconsideradas. 5.1 ESTATÍSTICAS BÁSICAS DA PROVA Esta seção apresenta estatísticas selecionadas e histogramas da nota geral ( 5.1.1 ) e de cada componente : Formação Geral ( 5.1.2 ) e Conhecimento Específico ( 5.1.3 ).
São também apresentadas estatísticas selecionadas de subpopulações, caracterizadas por grande região, categoria administrativa, organização acadêmica e modalidade de oferta. 5.1.1 ESTATÍSTICAS BÁSICAS GERAIS Na Tabela 5.1, são apresentadas as estatísticas básicas da prova, por grande região, dos estudantes concluintes de Engenharia Mecânica.
Como se pode verificar, a partir dos dados mostrados na Tabela 5.1, a média das notas da prova foi de 40,8 ( nas seções seguintes serão analisados os componentes de Formação Geral e de Conhecimento Específico ).
Os estudantes da região Norte obtiveram a média mais baixa ( 35,9 ) e os da região Nordeste alcançaram a média mais alta ( 41,3 ).
As demais médias foram 39,6, na região Centro-Oeste ; 41,0, na região Sudeste ; e 40,8, na região Sul.
O desvio-padrão para o Brasil foi de 12,9, sendo o maior encontrado na região Sudeste ( 13,2 ), e o menor, na região Norte ( 11,7 ).
A região que obteve a maior nota máxima foi a região Sudeste ( 91,6 ), ao passo que a região que atingiu a menor nota máxima foi o Norte ( 75,6 ).
A mediana do Brasil, como um todo, foi igual a 39,4, sendo a maior mediana obtida na região Nordeste ( 40,1 ) e a menor, na região Norte ( 34,6 ).
Considerando -se as notas segundo grande região, observa -se que existe diferença estatisticamente significativa ao nível de 95 \% entre as médias obtidas em todas as regiões, exceto entre Sudeste em relação ao Centro-Oeste e ao Nordeste, assim como entre o Sul em relação ao Centro-Oeste, ao Nordeste e ao Sudeste.
O comportamento das notas dos estudantes de todo o Brasil pode ser observado no Gráfico 5.1.
O Gráfico 5.1 apresenta um histograma com a distribuição das notas dos estudantes de Engenharia Mecânica de todo o Brasil : uma distribuição unimodal com a moda no intervalo ( 30 ; 40 ].
Na Tabela 5.2, são apresentadas informações referentes à média da nota final, desagregadas por categoria administrativa, organização acadêmica e modalidade de oferta dos cursos.
Os dados da Tabela 5.2 mostram que, em relação à categoria administrativa, os estudantes das IES públicas obtiveram média mais alta ( 47,8 ) do que a média das IES privadas ( 35,8 ).
O desvio-padrão para as IES públicas foi de 12,7, indicando uma dispersão das notas nesta categoria administrativa.
Já o desvio-padrão da IES privadas foi 10,5.
No que se refere à organização acadêmica, os estudantes de universidades obtiveram a média de 43,0, enquanto os de centros universitários, faculdades e CEFET/IF obtiveram, respectivamente, médias de 35,7, 35,2 e 44,4, respectivamente.
Constata -se que existe diferença estatisticamente significativa, ao nível de 95 \%, entre as médias das organizações acadêmicas.
A média da modalidade a distância ( 33,5 ) foi inferior à média da modalidade presencial ( 41,6 ), sendo a diferença entre as duas estatisticamente significativa, ao nível de 95 \%. 5.1.2 ESTATÍSTICAS BÁSICAS NO COMPONENTE DE FORMAÇÃO GERAL Na Tabela 5.3, são apresentadas as estatísticas básicas em relação ao componente da prova que avalia a Formação Geral dos estudantes concluintes.
A nota de Formação Geral considera os acertos nas nove questões objetivas e a nota média obtida na questão discursiva, sendo que a nota desta questão leva em conta 80 \% da nota relativa ao conteúdo e 20 \% da nota relativa ao desempenho linguístico. ( 47,4 ) e a maior, na região Nordeste ( 54,2 ).
A nota mínima nessa parte foi zero em quase todas as regiões, com exceção da região Centro-Oeste, em que a nota foi 7,2.
Considerando -se as notas, segundo grande região, observa -se que existe diferença estatisticamente significativa entre as médias das notas do componente de Formação Geral obtidas nas regiões, exceto entre as regiões Norte e Centro-Oeste ; Sudeste e Centro-Oeste ; Sudeste e Nordeste ; Sul e Centro-Oeste.
No Gráfico 5.2, é possível verificar a avaliação do desempenho dos estudantes no componente de Formação Geral, a partir do histograma da distribuição das notas correspondentes.
Nesse gráfico, a distribuição é unimodal, com moda em ( 50 ; 60 ].
Na Tabela 5.4, são apresentadas as informações referentes ao desempenho dos concluintes no componente de Formação Geral, em diferentes agregações : categoria administrativa, organização acadêmica e modalidade de oferta.
A partir dos dados da Tabela 5.4, observa -se que há diferença estatisticamente significativa entre as médias dos tipos de categoria administrativa.
A maior média foi obtida por estudantes de IES públicas ( 57,4 ) e a menor, pelos de IES privadas ( 47,8 ).
Em relação às médias por tipos de organização acadêmica, encontra -se : 53,5 para as universidades ; 47,9 para os centros universitários, 46,9 para as faculdades e 55,2 para os CEFET/IF.
A nota máxima obtida pelas IES privadas ( 99,7 ) foi menor em comparação às IES públicas.
O desvio-padrão menor foi percebido nas IES públicas ; e, entre as organizações acadêmicas, esse valor foi menor para os CEFET/IF.
Já no que se refere à modalidade de oferta, os cursos presenciais apresentaram média de 52,3, enquanto a média dos cursos a distância foi de 47,1.
A diferença é, portanto, estatisticamente significativa. 5.1.3 ESTATÍSTICAS BÁSICAS DO COMPONENTE DE CONHECIMENTO ESPECÍFICO Na Tabela 5.5, são apresentadas as estatísticas básicas referentes ao componente de Conhecimento Específico da área de Engenharia Mecânica.
A nota do componente de Conhecimento Específico leva em conta as notas da parte objetiva e a média das notas obtidas na questão da parte discursiva da área.
De acordo com a Tabela 5.5, a média do desempenho dos estudantes do Brasil para o componente de Conhecimento Específico da prova foi de 37,1.
A maior média foi obtida nas regiões Nordeste e Sul ( 37,4 ), e a menor, na região Norte ( 31,9 ).
As demais médias foram : 36,0, na região Centro-Oeste e 37,2, na região Sudeste.
Quanto à variabilidade das notas, o desvio-padrão do Brasil foi 14,5 ; sendo o maior desvio-padrão observado na região Sudeste ( 14,7 ) e o menor, na região Norte ( 12,8 ).
Os demais desvios foram : 14,1, na região Centro-Oeste e 14,3, nas regiões Nordeste e Sul.
A mediana das notas dos estudantes de todo o Brasil foi de 34,8.
A maior mediana ocorreu nas regiões Nordeste e Sul ( 36,8 ) e a menor, na região Norte ( 28,6 ).
As demais medianas foram : 32,8, na região Centro-Oeste e 34,8, na região Sudeste.
A nota máxima do Brasil, como um todo, foi 96,0 ; sendo obtida por, pelo menos, um estudante nas regiões Sudeste e Sul.
As demais notas máximas foram : 83,6, na região Centro-Oeste, 87,8 na região Nordeste e 79,6 na região Norte.
A nota mínima nessa parte foi zero para as regiões Nordeste e Sudeste, 8,2 para a região Centro-Oeste e 4,1 para as regiões Norte e Sul.
Observa -se que existe diferença estatisticamente significativa entre todas as regiões na análise das médias das notas do componente de Conhecimento Específico, exceto entre as regiões Nordeste e Centro-Oeste ; Sudeste em relação às regiões Centro-Oeste e Nordeste ; e Sul em relação às regiões Centro-Oeste e Sudeste.
Assim como nos Gráficos 5.1 e 5.2, mostrados anteriormente, no Gráfico 5.3, apresentado acima, foi possível ser realizada uma avaliação do desempenho de concluintes de Engenharia Mecânica, em relação ao componente de Conhecimento Específico, com um histograma da distribuição das notas correspondentes.
Esta também é uma distribuição unimodal, e o grupo modal é o ( 20 ; 30 ].
Na Tabela 5.6, é apresentada uma comparação dos resultados em relação à categoria administrativa, à organização acadêmica e à modalidade de oferta, levando -se em conta o desempenho de estudantes do componente de Conhecimento Específico da prova.
De acordo com os dados expostos na Tabela 5.6, no que se refere à organização acadêmica, a maior média foi dos CEFET/IF ( 40,8 ), seguida das universidades ( 39,4 ), dos centros universitários ( 31,6 ) e, por fim, das faculdades ( 31,3 ).
O maior desvio-padrão foi o das universidades ( 15,2 ).
Os CEFET/IF obtiveram a maior mediana ( 41,0 ).
Já as notas máximas foram : 96,0 nas universidades, 87,8 nos centros universitários, 83,8 nas faculdades e 91,8 nos CEFET/IF.
A nota mínima foi zero para as Universidades e Centros Universitários, e 4,1 para as faculdades e 1,0 para os CEFET/IF.
Observa -se que existe diferença estatisticamente significativa, ao nível de 95 \%, no componente de Conhecimento Específico entre as notas dos tipos de organização acadêmica.
Quanto à categoria administrativa, existe diferença estatisticamente significativa entre as médias das IES públicas ( 44,6 ) e as médias das IES privadas ( 31,8 ).
Neste caso, assim como no componente de Formação Geral, a maior média foi obtida por estudantes de IES públicas de ensino.
Os cursos de educação presencial apresentaram maior média ( 38,0 ) em relação aos cursos a distância ( 29,0 ), sendo a diferença estatisticamente significativa. 5.2 ANÁLISE DAS QUESTÕES OBJETIVAS Esta seção apresenta estatísticas selecionadas e histogramas das questões objetivas dos componentes de Formação Geral ( 5.2.1 ) e do componente de Conhecimento Específico ( 5.2.2 ).
São, também, apresentadas e comparadas as médias das subpopulações caracterizadas por grande região. 5.2.1 COMPONENTE DE FORMAÇÃO GERAL Na Tabela 5.7, são apresentadas as estatísticas básicas relativas às nove questões objetivas do componente da prova que abrange a Formação Geral dos estudantes.
De acordo com os dados da Tabela 5.7, em relação ao componente de Formação Geral, a média do Brasil foi 56,2.
A menor média foi encontrada na região Norte ( 50,6 ) e a maior, na região Nordeste ( 58,0 ).
As demais médias foram : 54,6, na região Centro-Oeste ; 56,6, na região Sudeste ; e 55,4, na região Sul.
O desvio-padrão do Brasil foi 19,4 ; sendo o maior desvio-padrão encontrado nas regiões Sudeste e Sul ( 19,4 ), e o menor desvio-padrão encontrado na região Norte ( 18,7 ).
A mediana do Brasil, como um todo, foi 55,6 a mesma encontrada em todas as regiões.
Quanto à nota máxima, todas as regiões alcançaram 100,0.
Assim como a nota mínima foi a mesma ( 0,0 ) para todas as regiões.
Na Tabela 5.8, exposta adiante, são apresentados o índice de facilidade e o índice de discriminação ( ponto-bisserial ) para cada uma das questões objetivas do componente de Formação Geral.
De acordo com os dados da Tabela 5.8, as questões objetivas do componente de Formação Geral, segundo o índice de facilidade, foram assim avaliadas : nenhuma teve o índice de facilidade classificado como “ Muito difícil ” ou como “ Muito fácil ” ; quatro questões foram classificadas como “ Fácil ” ; três questões foram classificadas com índice"Médio ” ; e duas questões foram classificadas na categoria “ Difícil ”.
O índice de facilidade variou de 0,26 a 0,79.
Já em relação ao índice de discriminação, seis questões foram classificadas com o índice “ Muito bom ” e três com o índice “ Bom ”.
Nenhuma questão foi classificada com o índice “ Médio ” ou"Fraco ”.
O índice de discriminação variou de 0,33 a 0,48.
Na Tabela 5.9, é apresentada a distribuição das questões, levando -se em conta, simultaneamente, a classificação dos dois índices.
Os dados da Tabela 5.9 atestam que seis questões que tiveram índice de discriminação “ Muito bom ” figuraram em dois níveis de dificuldade – “ Médio ” e “ Fácil ” : quatro classificadas na categoria “ Fácil ” ( questões 1, 2, 3 e 6 ) e duas classificadas na categoria “ Médio ” ( questões 4 e 8 ).
Em particular, as questões 3 e 6 foram as que apresentaram o maior poder discriminatório, com índice 0,48, considerado “ Fácil ” em termos de facilidade, com uma proporção de 0,65 e 0,77 de acertos, respectivamente, não tendo essas questões alcançado o máximo de acertos.
No Gráfico 5.4, para exemplificar, analisa -se o comportamento da Questão Objetiva 4 de Formação Geral.
Trata -se de uma questão considerada com nível “ Médio ” de dificuldade, em relação à facilidade, e que obteve o índice de discriminação 0,45 ( “ Muito bom ” ).
No Gráfico 5.4, cada uma das cinco curvas representa o percentual de respostas em determinada alternativa da questão, em função do número de acertos dos estudantes nessa parte da prova ( Formação Geral/múltipla escolha ), antes de possíveis eliminações pelo critério ponto-bisserial.
A curva em vermelho corresponde à alternativa A, a alternativa correta para essa questão.
Observa -se que, entre os estudantes com menor número de acertos, nessa parte do exame, a situação mais frequente foi a escolha da alternativa B ( em mostarda ), incorreta.
Por exemplo, entre os estudantes que acertaram até duas respostas, 14,4 \% escolheram a alternativa A ( em vermelho ) ; 35,0 \% escolheram a alternativa B ( em mostarda ) ; 7,8 \% escolheram a alternativa C ( em verde ) ; 10,1 \% escolheram a alternativa D ( em azul ) ; e 32,1 \% escolheram a alternativa E ( em lilás ).
Entre os que acertaram duas respostas entre as questões de múltipla escolha de Formação Geral, 0,7 \% deixou a questão em branco e ninguém marcou mais de uma alternativa, invalidando a questão.
À medida que o número de acertos aumenta, indicando desempenho melhor nessa parte da prova, aumenta, concomitantemente, a proporção de estudantes que selecionaram a alternativa correta A, atingindo 100 \% para os estudantes com oito acertos.
Essa análise permite verificar como a questão discriminou os grupos de desempenho, justificando -se o alto índice ( 0,45 ) obtido na questão.
Os gráficos relativos às demais questões de Formação Geral constam do Anexo IV. 5.2.2 COMPONENTE DE CONHECIMENTO ESPECÍFICO Na Tabela 5.10, são apresentadas as estatísticas básicas em relação às questões objetivas do componente de Conhecimento Específico da prova de Engenharia Mecânica por grande região.
Em relação ao componente de Conhecimento Específico da prova de Engenharia Mecânica, sete questões não foram usadas no cômputo das notas por terem sido descartadas pelo critério do ponto-bisserial.
Assim, as notas foram calculadas com base em 22 das 29 questões objetivas válidas do componente de Conhecimento Específico.
Como se verifica a partir dos dados da Tabela 5.10, a média do Brasil para o componente de Conhecimento Específico da prova foi 40,1.
A menor média foi observada na região Norte ( 34,8 ), e a maior, na região Sul ( 40,5 ).
O desvio-padrão de todo o Brasil foi 14,7, sendo o maior desvio-padrão encontrado na região Sudeste ( 15,0 ), e o menor, na região Norte ( 13,2 ).
A mediana de todo o Brasil foi 36,4, o mesmo valor da mediana encontrada nas regiões Centro-Oeste e Sudeste, e inferior às medianas das regiões Nordeste e Sul ( 40,9 ).
A nota máxima da prova foi 95,5, obtida, nas questões objetivas do componente de Conhecimento Específico, por, pelo menos, um estudante das regiões Sudeste e Sul, nas demais regiões, a nota máxima variou entre 77,3 ( Norte ) e 90,9 ( Centro-Oeste e Nordeste ).
A nota mínima de maior valor foi 9,1, obtida por, pelo menos, um estudante na região Centro-Oeste.
Na Tabela 5.11, são apresentados os índices de facilidade e de discriminação ( ponto-bisserial ) das questões objetivas do componente de Conhecimento Específico, para os estudantes de Engenharia Mecânica.
A partir do índice de facilidade obtido, apresentado na Tabela 5.11, pode -se concluir que a maioria das questões objetivas da prova foram classificadas no índice “ Difícil ” ( 19 ).
Cinco questões foram classificadas no índice"Fácil ” ; uma, no índice “ Médio ” ; uma, no índice “ Muito fácil ” ; e três questões foram alocadas no índice “ Muito difícil ”.
Quanto aos índices de discriminação das questões objetivas do componente de Conhecimento Específico da prova, obtém -se como resultado a seguinte classificação : quatro questões foram alocadas no índice “ Muito bom ” e seis, no “ Bom ”.
Dessa forma, dez das 29 questões foram classificadas com os índices de discriminação “ Bom ” e “ Muito bom ”.
Entre as demais, 12 delas foram alocadas no índice “ Médio ” e sete, no “ Fraco ”.
Constata -se, assim, que a prova — no que se refere ao componente de Conhecimento Específico — possuía média capacidade de discriminar entre estudantes que dominam ou não o conteúdo.
O índice de facilidade variou de 0,10 a 0,90, e o de discriminação, de 0,10 a 0,55.
Na Tabela 5.12, é apresentada a distribuição das questões considerando -se, simultaneamente, a classificação dos dois índices.
Os dados da Tabela 5.12 atestam que seis questões que tiveram índice de discriminação “ Bom ” figuraram em dois níveis de dificuldade — “ Difícil"e"Fácil ” : quatro classificada na categoria “ Difícil”e duas, na categoria “ Fácil ”.
As questões que tiveram o índice de discriminação “ Fraco ” configuraram -se com índices de facilidade “ Médio ” ( uma questão ), “ Difícil ” ( quatro questões ) e “ Muito Difícil ” ( duas questões ).
Entre as questões que alcançaram os maiores índices de discriminação, “ Muito bom ”, uma delas, a questão 15, teve o maior índice de discriminação ( 0,55 ).
Tal questão foi classificada como “ Difícil ”, com índice de facilidade igual a 0,38.
A questão 22 foi a mais “ Difícil ” entre as 29 questões específicas válidas, com baixo índice de facilidade, apenas 0,10 de proporção de acertos.
Essa questão apresentou poder discriminatório fraco, 0,11, o que comprova ter sido uma das mais difíceis para os estudantes.
Destacam -se, também, com baixo índice de facilidade, as questões 21, 33 e 22, cujos índices de facilidade, em termos proporcionais, correspondem, respectivamente, a 0,15, 0,11 e 0,10 de estudantes que responderam acertadamente.
Essas três questões tiveram poder discriminatório “ Médio","Fraco"e"Fraco ”, com índices 0,20, 0,11 e 0,13, respectivamente.
Destaca -se que as questões 10, 14, 22, 29, 32, 33 e 37 foram desconsideradas pelo critério ponto-bisserial.
Por isso, foram eliminadas do cômputo da nota final.
A título de exemplo das análises do comportamento das questões objetivas, no Gráfico 5.5, analisa -se a Questão 17 do componente de Conhecimento Específico.
No Gráfico 5.5, cada uma das cinco curvas representa o percentual de respostas em determinada alternativa da questão, em função do número de acertos dos estudantes nessa parte da prova ( Componente Específico ), antes de possíveis eliminações pelo critério ponto-bisserial.
A alternativa correta A, representada no gráfico pela curva em vermelho, foi escolhida em maiores proporções pelos estudantes com desempenho melhor nessa parte da prova.
Já as alternativas incorretas, também denominadas distratores, foram selecionadas, principalmente, por aqueles com notas mais baixas.
A questão 17 foi considerada, pelas respostas dos estudantes avaliados na prova, como uma questão com índice de facilidade “ Fácil, 0,67 ”, ou seja, 67,0 \% dos estudantes assinalaram, acertadamente, a opção A correspondente ao gabarito.
Seu índice de discriminação foi igual a 0,42, classificado como “ Muito bom ”.
Observa -se que a soma não é 100 \%, em razão das questões não respondidas ou com mais de uma opção marcada.
Cerca de 100,0 \% daqueles com nota zero deixaram a questão em branco ou fizeram dupla marcação.
A proporção de estudantes que selecionaram a resposta correta A aumenta gradativamente, chegando a atingir 100 \% para 19 acertos, enquanto a proporção dos que escolheram alternativas incorretas decai, a partir de seis acertos, em função do número de acertos nessa parte da prova.
Os gráficos relativos às demais questões do componente de Conhecimento Específico constam no Anexo IV. 5.3 ANÁLISE DAS QUESTÕES DISCURSIVAS Esta seção apresenta estatísticas selecionadas e histogramas das questões discursivas de Formação Geral ( 5.3.1 ) e de Conhecimento Específico ( 5.3.2 ).
Expõe e compara, também, as médias de subpopulações, caracterizadas por grande região.
Cumpre notar que nem todas as questões passam pelo crivo dos corretores.
Após a digitalização, as respostas passam por um pré-processamento de identificação automático.
Algumas dessas questões são definidas como um “ branco automático ” ( que não são enviadas para os professores para correção, a não ser que, na segunda fase, a de inspeção visual, seja constatada uma resposta ).
O critério usado nesse pré-processamento consiste em verificar, para cada questão, a partir da quantidade de pixels, o que, provavelmente, é uma resposta em branco.
A partir da separação, os documentos que foram considerados em branco são enviados para uma verificação visual, feita por uma equipe de colaboradores devidamente capacitados.
Desse modo, qualquer erro, presente no pré-processamento, é encontrado, e as questões preenchidas são enviadas para a correção dos avaliadores.
Eventualmente, algumas respostas em branco não são detectadas nesse pré-processamento automático.
A presença de pequenos riscos, ou manchas, pode fazer o programa não identificar uma questão como um possível branco.
Essas são enviadas para a correção, mas são classificadas como “ branco ” pelos avaliadores, no momento da correção. 5.3.1 COMPONENTE DE FORMAÇÃO GERAL As análises dos resultados de desempenho dos estudantes de Engenharia Mecânica, na questão discursiva relativa à Formação Geral, encontram -se na Tabela 5.13 e no Gráfico 5.6, cujos dados são analisados de forma sequenciada.
No caso da Tabela 5.13, analisa -se a média das notas da questão discursiva de Formação Geral, que leva em conta as correções realizadas por duas bancas : uma que analisa e pontua aspectos relativos ao conteúdo das respostas, seguindo o padrão de respostas divulgado pelo Inep ( ver Anexo VIII ) ; e outra que avalia o desempenho linguístico expresso pela redação das respostas, seguindo um padrão de respostas específico da área de Língua Portuguesa, detalhado na seção 5.3.1.5 deste relatório.
No cálculo da nota da questão discursiva, consideram -se 80 \% relativos à nota da correção de conteúdo e 20 \% referentes à nota de desempenho linguístico.
No Gráfico 5.6, está representada a distribuição das notas na questão discursiva do componente de Formação Geral.
Na Tabela 5.13, observa -se que os estudantes de todo o Brasil obtiveram, em Formação Geral, média de 43,6 nas questões discursivas.
Pode -se notar, também, que o desvio-padrão nesse conjunto de questões foi de 24,7.
A maior média foi obtida na região Sudeste ( 44,4 ), e a menor, na região Sul ( 42,5 ).
A mediana de todo o Brasil, neste componente, foi 44,5, e na região Sudeste ( 45,4 ) foi encontrada a maior mediana.
A nota máxima ( 100,0 ) foi obtida por, pelo menos, um estudante das regiões Nordeste, Sudeste e Sul.
A nota mínima ( 0,0 ) foi a mesma em todas as regiões do Brasil.
No Gráfico 5.6, está representada a distribuição das notas na questão discursiva do componente de Formação Geral.
Como mostra o Gráfico 5.6, a moda da distribuição ocorre no intervalo ( 50 ; 60 ], com frequência de 15,3 \% dos participantes, sendo que, no intervalo [ 0 ; 10 ], inclui -se, além da nota zero, a frequência de estudantes que deixaram esse tipo de questão em branco.
Na sequência, os resultados verificados para a questão discursiva de Formação Geral são apresentados, estabelecendo -se relações com os conteúdos abordados.
Os comentários da banca de docentes avaliadores, a respeito do observado na correção das respostas dos estudantes, suas impressões e conclusões serão apresentadas junto à análise da questão.
Cumpre esclarecer que, tendo em vista que a questão discursiva de Formação Geral foi padronizada, ou seja, foi a mesma em todas as provas, os comentários da banca são os mesmos para todas as carreiras acadêmicas, sendo direcionados a todos os estudantes que participaram do Enade 2023.
A seguir, serão analisados o desempenho linguístico dos estudantes da área de Engenharia Mecânica na questão discursiva de Formação Geral do Enade 2023, apresentando -se os resultados obtidos e os comentários da banca de correção sobre a questão. 5.3.1.1 ANÁLISE DE CONTEÚDO DA QUESTÃO DISCURSIVA DO COMPONENTE DE FORMAÇÃO GERAL Os dados de Engenharia Mecânica, obtidos a partir das respostas à questão discursiva do componente de Formação Geral, encontram -se na Tabela 5.14 e no Gráfico 5.7, exposto após a apresentação dos dados da referida tabela.
Verifica -se, a partir dos dados da Tabela 5.14, que, no conteúdo da questão discursiva do componente de Formação Geral, os estudantes de todo o Brasil obtiveram média de 38,9.
A maior média foi obtida na região Sudeste ( 39,6 ), e a menor, na região Sul ( 37,7 ).
Quanto à variabilidade das notas, o desvio-padrão de todo o Brasil foi 27,9.
O menor desvio-padrão foi obtido na região Sul ( 27,7 ), e o maior desvio-padrão foi obtido pela região Norte ( 28,5 ).
A mediana do Brasil, como um todo, foi 37,5, sendo igual à das regiões Centro-Oeste, Nordeste, Norte e Sul ( 37,5 ) e inferior à mediana da região Sudeste ( 39,0 ).
A nota máxima da questão discursiva foi 100,0, obtida por, pelo menos, um estudante em todas as regiões.
Já a nota mínima para todas as regiões foi zero.
No Gráfico 5.7, apresenta -se a distribuição das notas de conteúdo da questão discursiva do componente de Formação Geral.
O percentual de participantes que deixaram em branco a resposta ou zeraram a questão discursiva na avaliação de conteúdo foi 17,6 \%.
O Gráfico 5.7 mostra que a moda dessa distribuição ocorre no intervalo [ 0 ; 10 ], com uma frequência de 21,4 \% dos participantes. 5.3.1.2 COMENTÁRIOS SOBRE A CORREÇÃO DE CONTEÚDO DAS RESPOSTAS À QUESTÃO DISCURSIVA A questão discursiva do componente de Formação Geral foi composta por um texto motivador e por dois itens ( a e b ).
O texto motivador apresentou um comentário sobre uma publicação do Instituto Brasileiro de Geografia e Estatística ( IBGE ) que tratava dos indicadores sobre a realidade social brasileira, apontando o crescimento do nível de desigualdade socioeconômica em 2021 e concluindo que a pobreza teve aumento recorde nesse período.
A questão foi dividida em dois itens.
O primeiro item solicitava que o respondente estabelecesse uma relação entre o perfil da população mais afetada pelas desigualdades sociais e as situações de risco socioambiental.
Já o segundo item pedia a proposição de duas possíveis ações a serem desenvolvidas em bairros periféricos a fim de minimizar os impactos socioambientais.
A questão foi pertinente, apresentou linguagem acessível e abordou conteúdo basilar para Formação Geral ; no entanto, apresentou um nível de dificuldade alto, dadas as competências intrincadas que demandou.
O tema foi atual e de grande relevância para a sociedade brasileira.
Abordou assunto de impacto social e de provável domínio dos concluintes de curso superior.
O texto motivador foi adequado e coerente com o que se pediu na questão, tratou do aumento da desigualdade social e da descrição do perfil da população com base em critérios socioeconômicos.
Com relação às habilidades cognitivas, observou -se que as seguintes foram necessárias para responder à questão proposta : capacidade de expressar ideias de maneira clara, coesa e argumentativa por meio da produção de textos ; capacidade de inter-relacionar objetos de conhecimento nas diferentes áreas ; e capacidade de criar estratégias de ação, selecionar métodos e formular propostas.
Entre os saberes necessários para o desenvolvimento das respostas, os respondentes deveriam conhecer o meio social e/ou geográfico da população brasileira de perfil socioeconômico pobre ou extremamente pobre e deveriam saber relacionar esse meio com situações de risco socioambiental que afligem principalmente esse perfil populacional.
Portanto, deveriam, também, conhecer uma variedade de riscos socioambientais a que esse perfil populacional está sujeito, por conta do meio social e/ou geográfico em que vive.
A questão, em seu item b ), por sua vez, demandou o domínio de um rol de propostas minimamente factíveis para reduzir tais situações de risco, envolvendo distintos atores sociais.
Em relação ao padrão de resposta, esperava -se para cada um dos itens avaliados que : a ) O estudante explicasse a relação entre desigualdade social e riscos socioambientais percorrendo o seguinte trajeto teórico-argumentativo :  descrever o(s) risco(s) socioambiental(is) que afeta(m) a população com perfil pobre ou extremamente pobre, tais como inundações, deslizamentos de terra, contaminação ambiental, maior vulnerabilidade em relação a doenças, discriminação social e criminalidade, entre outros ;  descrever o meio ( geográfico ou social ), como morros, áreas próximas a rios, mangues, espaços marcados pela violência urbana, entre outros, nas cidades brasileiras, como fator de criação ou potencialização das condições de vulnerabilidade e risco aos pobres ou extremamente pobres ;  estabelecer nexo causal entre a pobreza ou extrema pobreza, o meio ( geográfico ou social ) urbano e a situação de risco socioambiental, relacionando esses aspectos ao perfil socioeconômico da população.
Portanto, esperava -se que o estudante apresentasse consistência argumentativa quanto ao que foi solicitado, notadamente quanto aos riscos socioambientais, à descrição do meio e ao estabelecimento do nexo causal entre pobreza, meio e risco socioambiental. b ) O estudante apresentasse propostas pertinentes, factíveis e bem desenvolvidas que envolvessem ação governamental ( federal, estadual/distrital ou municipal ) e participação da comunidade a fim de minimizar riscos socioambientais, como, por exemplo :  implantar obras de infraestrutura urbana que envolvessem contenção de morros e encostas ; saneamento básico : serviços regulares de abastecimento de água, coleta e tratamento de esgotos, coleta e manejo de resíduos sólidos, manejo de águas pluviais, entre outras ;  desenvolver políticas sociais, como segurança pública, provisão de moradia adequada, com custo acessível, e regularizada do ponto de vista fundiário ; educação de qualidade ; atendimento psicossocial, com especial atenção para grupos em situação de vulnerabilidade ; entre outras ;  garantir o acesso dessas comunidades a equipamentos sociais e culturais, bem como a espaços públicos inclusivos e a áreas verdes ;  elaborar planos estratégicos participativos urbanos e ambientais a serem desenvolvidos nas comunidades em situação de vulnerabilidade.
No tocante ao subitem b ), esperava -se que o respondente apresentasse coerência propositiva, englobando os diferentes atores, quais sejam : o governo, em suas diversas instâncias, e a comunidade, também em suas diversas organizações, almejando conter os riscos socioambientais. 5.3.1.3 CONSIDERAÇÕES FINAIS SOBRE A CORREÇÃO DE CONTEÚDO DA QUESTÃO DISCURSIVA DO COMPONENTE DE FORMAÇÃO GERAL De uma maneira geral, os respondentes apresentaram um nível aquém do esperado em termos de formação geral, destacando -se, neste âmbito, os aspectos a seguir.
Nas respostas relacionadas ao item a ), percebeu -se que, quanto aos riscos, a maioria dos respondentes não se limitou a descrevê -los, ou mesmo a estabelecer um rol deles.
No caso daqueles que descreveram as diversas situações de risco, em sua maioria, as respostas se enquadraram naquilo que estava previsto no padrão de resposta, apresentando os riscos derivados da insegurança relativa ao meio geográfico em que habitam, tais como alagamentos, deslizamentos e doenças por contato com um meio insalubre.
Quanto à descrição do meio, notou -se que houve uma maior aderência por parte dos respondentes, pois a maior parte das respostas, excluindo -se as que foram deixadas em branco, descreveu o meio de alguma forma.
Devido à grande variedade de possíveis descrições do meio geográfico e/ou social em que vive a população brasileira de perfil pobre ou extremamente pobre, o padrão de respostas deixava aberta a possibilidade de considerar respostas variadas em relação às citadas como exemplo em sua redação.
E, de fato, entre as respostas encontradas, verificaram -se muitas que foram consideradas certas, por estarem dentro do perfil previsto como correto, mas que eram distintas dos exemplos específicos apresentados no padrão de respostas.
O mesmo se deu quanto ao item b ) : houve uma grande variedade de proposições para além dos campos temáticos específicos previstos no padrão de respostas, mas que foram consideradas corretas porque atendiam à demanda do item.
Para esse item, o padrão de respostas demandava que as Dessa forma, o desempenho dos estudantes foi, em geral, muito baixo.
O número de questões deixadas em branco foi considerável, bem como o quantitativo de respostas muito curtas, ocupando menos de cinco das 15 linhas oferecidas para a resposta.
Muitos responderam a apenas um dos dois itens.
A grande maioria dos que buscaram responder ao que a questão demandava, completamente ou parcialmente, não souberam se expressar de modo correto, objetivo e coerente.
Não poucos respondentes repetiram longos trechos do enunciado em suas respostas, demonstrando um baixo nível de formação geral e, também, baixa disposição para a realização da prova.
Muitos respondentes foram desatentos e escreveram a resposta à questão discursiva de componente de Conhecimento Específico no campo dedicado à questão discursiva de Formação Geral, tendo sido incluídos dentre os respondentes que fugiram ao tema.
Uma quantidade proporcionalmente pequena, mas grande numericamente, foi a de estudantes que apresentaram respostas que desrespeitavam os direitos humanos, em sua maioria propondo o controle reprodutivo sobre a população pobre ou extremamente pobre como solução governamental para as situações de risco socioambiental comumente vivenciadas pelas populações brasileiras pobres ou extremamente pobres.
Alguns estudantes também apontaram uma relação necessária entre pobreza e ignorância ou mesmo entre pobreza e violência.
Entre os estudantes que responderam com maior dedicação, buscando dar conta de toda a questão, houve um número significativo que aparentou ter grandes dificuldade de expressão textual, dando a entender que seu baixo domínio da linguagem verbal comprometia a capacidade de produzir respostas adequadas.
Muitas respostas não apresentaram o mínimo de lógica ou não buscaram estabelecer relações entre os conteúdos mobilizados.
Portanto, em sua maioria, os estudantes não apresentaram a capacidade de expressar ideias de maneira clara, coesa e argumentativa por meio da produção de textos, uma habilidade fundamental ao concluinte do Ensino Superior.
Um outro perfil de destaque foi o de estudantes que se dedicaram a pontuar elementos válidos para responder à questão sem utilizar os conectivos, as relações causais.
Também merecem destaques os respondentes que falaram sobre aspectos históricos da pobreza e a descrição dos tipos sociais mais envolvidos em situação de pobreza ou extrema pobreza, fugindo de buscar estabelecer uma relação causal com as situações de risco socioambiental, portanto, fugindo do que a questão solicitava.
Essas situações demonstraram um perfil majoritário com baixa capacidade de produzir para o item de Formação Geral uma resposta adequada esperada para estudantes concluintes do Ensino Superior.
Os estudantes com perfis de respostas minimamente adequadas ou boas corresponderam ao segundo perfil majoritário.
Em sua maior parte, esse perfil de estudantes respondeu bem ao subitem a ), mas não respondeu bem ou subitem b ), ou vice-versa.
O perfil minoritário foi o de estudantes de nível muito bom a ótimo.
Esses com competência para descrever objetivamente os pontos essenciais da questão, estabelecendo o nexo causal entre eles.
Quanto à relação com o conteúdo, a maior parte dos respondentes soube pensar nos elementos e nas imagens do problema levantado na questão, mas não soube fazer as ligações que estimulam a relação entre esses elementos.
Portanto, demonstraram uma baixa capacidade de inter-relacionar objetos de conhecimento nas diferentes áreas, outra habilidade esperada de um concluinte do Ensino Superior.
Uma grande parte dos estudantes vê o saneamento básico precário como a maior fonte de riscos socioambientais para a população pobre ou extremamente pobre.
As descrições do meio físico, como as situações de moradias em encostas de morros e beiras de rio, também constituíram entre as mais comuns.
Em menor número, mas ainda com relativa frequência, apareceram os casos de estudantes que apontaram o meio como lugar de domínio territorial de criminosos ou os casos que indicaram o meio social como opressivo aos negros.
Quanto às proposições, os estudantes, em sua maioria, não souberam criar estratégias de ação, selecionar métodos e formular propostas.
Uma quantidade considerável de estudantes apresentou proposições muito genéricas como “ construir escolas ”.
A maioria não apontou conexão entre as propostas e a minimização dos riscos socioambientais, como se a ação proposta já evidenciasse o que ela está solucionando, assim também indicando uma grande concentração de estudantes que não aparenta dominar habilidades importantes para concluintes do Ensino Superior. 5.3.1.4 ANÁLISE DE LÍNGUA PORTUGUESA DA QUESTÃO DISCURSIVA DO COMPONENTE DE FORMAÇÃO GERAL Os dados de Engenharia Mecânica obtidos a partir das respostas à questão discursiva do componente de Formação Geral, no que tange ao desempenho linguístico, são apresentados na Tabela 5.15.
Os dados da Tabela 5.15 mostram que, no quesito desempenho linguístico, os estudantes, de todo o Brasil, obtiveram média de 62,6.
A maior média em relação ao desempenho linguístico foi obtida na região Sudeste ( 63,3 ), e a menor, na região Centro-Oeste ( 60,9 ).
Quanto à variabilidade das notas, o desvio-padrão de todo o Brasil foi 23,4.
O menor desvio-padrão foi obtido na região Sudeste ( 22,9 ), e o maior desvio-padrão foi obtido na região Nordeste ( 24,4 ).
A mediana das notas de Língua Portuguesa foi 67,5 para o Brasil, como um todo, a mesma obtida em todas as regiões.
A nota máxima para todo o Brasil foi 100,0, com, pelo menos, um estudante com essa nota em todas as regiões.
A nota mínima foi 0,0 em todas as regiões do país.
No Gráfico 5.8, mostra -se a distribuição das notas de desempenho linguístico do componente de Formação Geral.
O Gráfico 5.8 demonstra que 8,1 \% dos participantes correspondem ao intervalo [ 0 ; 10 ].
A moda dessa distribuição ocorre no intervalo ( 70 ; 80 ], com uma frequência de 26,6 \% dos participantes. 5.3.1.5 COMENTÁRIOS SOBRE A CORREÇÃO DE LÍNGUA PORTUGUESA DAS RESPOSTAS À QUESTÃO DISCURSIVA Os comentários feitos neste item a respeito da correção de Língua Portuguesa das respostas emitidas para a questão discursiva do componente de Formação Geral estão agrupados em três pontos : avaliação da questão, padrão de resposta utilizado na avaliação e análise das respostas dadas pelos estudantes.
A linguagem da questão e do texto motivador foi considerada condizente com o domínio linguístico que se espera de estudantes concluintes de cursos em nível de graduação, sendo -lhes, portanto, acessível.
Do ponto de vista linguístico e textual, o nível da questão foi considerado mediano.
Por meio de um comando claro e objetivo, a questão solicitava que o estudante explicasse determinada relação, exposta no texto motivador, e apresentasse duas propostas para atenuar o problema discutido.
Para compreender o que deveria ser feito, o estudante necessitava ter domínio da leitura e interpretação textual, de modo a entender o enunciado, estabelecer a relação proposta na questão e apontar possíveis ações governamentais e da comunidade para resolver os problemas elencados.
O fato de a proposta textual ser apresentada em dois itens ( a e b ) separadamente levou muitos estudantes a produzirem o texto em dois parágrafos bem definidos e, na maioria dos casos, sem conexão textual entre ambos, ou seja, respondendo à questão em dois ou mais parágrafos que indicavam a letra do item ao qual a resposta se dirigia.
Foram relativamente poucas as provas em que os estudantes respondiam à questão em um texto que articulasse textualmente os dois itens.
Essa formatação dificultou a avaliação do quesito Aspectos Textuais no que se refere à organização textual.
Como a questão demandava certa capacidade de análise, a resposta exigia, em termos textuais, que o estudante fosse além do texto motivador, o que, em parte, passou a identificar a qualidade textual da resposta, ainda que tomasse como base o próprio enunciado da questão e empregasse dados do texto motivador.
Textos com qualidade mediana, em geral, consistiram na retextualização da questão e/ou de enunciados do texto motivador, com alguma autonomia linguística para elaborar argumentação superficial.
A simples cópia — às vezes, com erros ortográficos ou morfossintáticos — da questão ou de enunciados do texto motivador, praticamente sem autonomia textual ou argumentativa, foi um recurso característico de respostas de baixa qualidade em termos linguísticos, porém não penalizado na avaliação do desempenho linguístico.
Em relação ao padrão de resposta utilizado na avaliação, foram considerados os seguintes aspectos : ( 1 ) aspectos ortográficos ; ( 2 ) aspectos textuais ; ( 3 ) aspectos morfossintáticos e de propriedade vocabular adequada à modalidade escrita formal da Língua Portuguesa, exigida pela situação comunicativa.
A seguir, são apresentados cada um dos aspectos avaliados : 1 ) Aspectos ortográficos – referem -se ao domínio das convenções ortográficas da norma-padrão da Língua Portuguesa, como : grafia de vogais e consoantes, uso de maiúsculas e minúsculas, emprego do hífen e da acentuação gráfica, segmentação de palavras. 2 ) Aspectos textuais – referem -se ao domínio dos procedimentos de estruturação textual : organização interna dos períodos, emprego de conectores para a articulação lógica entre os períodos e entre os parágrafos, emprego de marcas de referenciação lexical e pronominal ; e utilização dos sinais de pontuação que contribuem para a organização lógica da frase. 3 ) Aspectos morfossintáticos e de propriedade vocabular – referem -se ao domínio das regras de caráter morfossintático estabelecidas como modelares do ponto de vista da modalidade escrita formal da norma-padrão da Língua Portuguesa : concordância nominal e verbal, regência nominal e verbal, colocação pronominal, flexão nominal e verbal, correlação entre tempos e modos verbais, ausência de marcas de oralidade.
Outra competência combinada a esta na análise das respostas foi a seleção vocabular adequada à modalidade escrita formal da Língua Portuguesa, tendo em vista a interseção entre as duas do ponto de vista das exigências do registro formal da modalidade escrita da norma-padrão.
No que concerne à análise das respostas dadas pelos estudantes, especificamente, em relação ao desempenho linguístico, foram observados muitos deslizes gramaticais no que diz respeito à acentuação e grafia correta das palavras, embora, de maneira geral, os estudantes tenham demonstrado domínio da ortografia vigente.
Quanto aos aspectos textuais, verificou -se predominância de textos organizados sem a presença das três partes esperadas para um texto dissertativo-argumentativo : introdução ; desenvolvimento e conclusão.
A argumentação foi predominantemente rasa e pouco elaborada, havendo frequentes desvios de pontuação, nomeadamente quanto ao emprego da vírgula.
Por fim, quanto aos aspectos morfossintáticos, observou -se um bom domínio das regras de regência e concordância na maior parte dos textos.
A adequação lexical também foi pertinente a uma redação formal na maior parte das vezes.
Em relação ao item a ), destacaram -se como respostas a relação entre comunidades negras periféricas, além de ribeirinhos e indivíduos em vulnerabilidade socioeconômica como sendo o perfil da população brasileira mais propenso a riscos socioambientais.
No que diz respeito ao item b ), foram encontradas com maior frequência propostas relacionadas a atuações governamentais, familiares e sociais diversas, como a necessidade de maiores investimentos e ofertas de trabalho na periferia.
Houve a ocorrência de respostas fracas e incompletas.
O item a ) foi respondido, na maior parte dos textos, de modo discursivo, enquanto o item b ) figurou, por vezes, como uma lista de ações que poderiam resolver o problema.
Em alguns casos, não tão raros, foram produzidos textos nos quais os estudantes se recusaram explicitamente a produzir uma resposta adequada, alegando que a questão era pertinente a uma área de formação diferente da sua.
As respostas encontradas e mesmo a ausência de respostas nos permitem tecer algumas reflexões.
Há constante falta de atenção na leitura do enunciado da questão, o que levou alguns estudantes a responderem à questão específica em local diverso do esperado, prejudicando a resposta.
Por vezes, alguns candidatos produziram as respostas para as duas questões ( Formação Geral e Conhecimento Específico ) no mesmo espaço, prejudicando o desempenho textual e argumentativo.
Seja por desatenção, seja por descuido, a dificuldade de interpretação é sintomática e deve ser considerada.
Desvios ortográficos, morfossintáticos e de pontuação frequentes apontam para lacunas no que concerne à alfabetização e ao desenvolvimento da escrita ainda nos anos do ensino fundamental.
Contudo, o número baixo de textos que tenham obtido notas muito baixas, como 1,0 e 2,0, revela que, entre os estudantes avaliados, não havia pessoas com níveis extremamente baixos de alfabetização e letramento.
Problemas quanto à organização do texto dissertativo-argumentativo, com ausência de introdução, desenvolvimento e conclusão, mesmo que minimamente, apontam para a necessidade de um trabalho mais profundo com a escrita em cursos superiores.
A argumentação rasa ou mesmo o desinteresse por temas sociais revelam o fato de que alguns dos cursos estudados não trabalham o senso crítico do estudante, não o conduzindo a uma vivência plenamente crítica e democrática em sociedade.
A dificuldade em responder ao item b ) da questão demonstrou uma fragilidade teórica na formação, principalmente no que concerne à atuação do profissional em relação a questões de cunho socioeconômico e ao atendimento ao público mais fragilizado socioeconomicamente.
A baixa frequência de desvios morfossintáticos pode ser justificada pela própria competência linguística do estudante, demonstrando que um trabalho mínimo com questões gramaticais nos cursos superiores poderia sanar grande parte do problema.
Relativamente poucos estudantes demonstraram incapacidade quanto à escolha de palavras adequadas para a produção de um texto formal, apresentando competência genérica, na maior parte das vezes, satisfatória.
Textos que obtiveram nota máxima foram raros, o que aponta para a necessidade de reforçar a importância do trabalho com a leitura e escrita em âmbito formal com os estudantes dos cursos avaliados. 5.3.1.6 CONSIDERAÇÕES FINAIS SOBRE A CORREÇÃO DE LÍNGUA PORTUGUESA DAS RESPOSTAS À QUESTÃO DISCURSIVA DO COMPONENTE DE FORMAÇÃO GERAL A análise global em relação ao desempenho linguístico dos cursos avaliados no Enade 2023 revelou um desempenho mediano dos estudantes avaliados.
Quanto aos aspectos textuais e argumentativos, verificou -se um nível baixo de conhecimento e mesmo de interesse, por parte de muitos estudantes, em refletir sobre os problemas sociais envolvendo pessoas em situação de vulnerabilidade e fenômenos de risco socioambiental.
A maioria dos estudantes não conseguiu apresentar uma análise crítica, tampouco responder conforme solicitado pela questão.
Alguns estudantes não foram atentos às informações apresentadas nos textos nem ao que os enunciados solicitavam.
Alguns demonstraram o domínio sobre o assunto ou, no mínimo, o comprometimento e o interesse para responder as questões.
Entende -se que o tempo, a quantidade de questões e a pressão podem influenciar no momento de elaborar a resposta.
Além disso, alguns estudantes de cursos específicos não se sentiram motivados a falar de um assunto julgado como “ social ” e não pertinente à sua área de atuação profissional, o que é um forte representativo de vieses ideológicos que perpassam determinadas formações.
Cabe ressaltar que a questão abordava um conteúdo atual e relevante a ser considerado por um indivíduo plenamente capaz de agir em sociedade de forma crítica, de qualquer área de formação em nível superior.
Destaca -se que estudantes com aparente domínio da área de saúde apresentaram inclinação crítica a abordar o tema de forma bastante humanizada.
Quanto aos aspectos ortográficos e gramaticais de natureza morfossintática, observam -se muitas lacunas trazidas desde a educação básica.
Entende -se que a universidade não pode sanar todos os problemas de cunho ortográfico e gramatical, mas é preciso salientar a importância de disciplinas voltadas à leitura e produção de texto serem incluídas como disciplinas obrigatórias nos cursos avaliados. 5.3.2 COMPONENTE DE CONHECIMENTO ESPECÍFICO As análises dos resultados de desempenho dos estudantes de Engenharia Mecânica, na questão discursiva relativa ao componente de Conhecimento Específico, encontram -se na Tabela 5.16 e no Gráfico 5.9. 5.3.2.1 ANÁLISE DA QUESTÃO DISCURSIVA DO COMPONENTE DE CONHECIMENTO ESPECÍFICO A Tabela 5.16 apresenta os dados sobre a questão discursiva do componente de Conhecimento Específico.
A partir dos dados registrados na Tabela 5.16, observa -se que a média foi mais baixa se comparada à média obtida na questão discursiva do componente de Formação Geral.
Enquanto, no componente de Formação Geral, a média para estudantes de Engenharia Mecânica de todo o Brasil foi 43,6 ; na parte de Conhecimento Específico, a média foi 9,5.
A maior média desse componente foi obtida pelos estudantes das regiões Centro-Oeste e Nordeste ( 10,0 ), e a menor, pelos da região Norte ( 5,8 ).
Quanto à variabilidade das notas, o desvio-padrão de todo o Brasil foi 23,7.
O maior desvio-padrão foi encontrado na região Centro-Oeste ( 24,9 ), e o menor, na região Norte ( 20,3 ).
A nota máxima obtida foi 100,0 em todas as regiões, assim como a nota mínima ( 0,0 ), também, foi a mesma em todas as regiões do país.
No Gráfico 5.9, apresenta -se a distribuição das notas na questão discursiva no componente de Conhecimento Específico.
À luz dos dados do Gráfico 5.9, verifica -se que a moda da distribuição das notas da questão discursiva do componente de Conhecimento Específico ocorre no intervalo [ 0 ; 10 ], com 83,8 \% do total de participantes, um valor considerável.
Vale ressaltar que, nesse intervalo, são contabilizadas também as questões deixadas em branco. 5.3.2.2 COMENTÁRIOS SOBRE AS RESPOSTAS À QUESTÃO DISCURSIVA A questão discursiva aplicada para a área de Engenharia Mecânica foi composta por um texto motivador seguido de uma breve descrição de um sistema mecânico apresentado em forma de desenho esquemático, com a descrição das medidas necessárias para a resolução da questão.
Os quesitos propostos para a análise da questão foram apresentados de forma clara, separados em dois itens ( a e b ).
No item a ), foram solicitados os valores dos módulos dos torques presentes em três trechos do sistema de transmissão – AB, BC e CD.
Já no item b ) foram solicitados os valores das reações apresentadas nos dois apoios do sistema.
O tema proposto na questão foi pertinente aos conhecimentos adquiridos no curso de engenharia, abordando conteúdos fundamentais para a formação do engenheiro mecânico.
A questão apresentou uma dificuldade média em relação aos conceitos de análise de tensões atuantes em sistemas mecânicos de transmissão.
O padrão de resposta para a questão foi desenvolvido de forma assertiva sem possibilidades de resposta diversa da apresentada, sendo imprescindível a informação correta tanto das reações quanto dos seus valores, não possibilitando correções parciais.
Assim quaisquer textos explicativos sem a apresentação dos valores corretos foram avaliados binariamente com nota zero para cada um dos padrões apresentados.
Em relação às respostas dos estudantes, o padrão da maioria das respostas foi apresentado com a questão totalmente em branco ou com alguns desenhos esquemáticos, mas sem o desenvolvimento dos cálculos.
Somando -se as questões em branco, outras respostas se limitaram a explicações sobre esforços de torção ou utilizaram valores errados para os cálculos.
Uma parcela dos estudantes foi capaz de resolver o item a ), mas não o item b ), o que pode indicar dificuldade na análise de esforços em diferentes planos.
Outra parcela de estudantes apresentou desenvolvimento correto para o item b ), mas errou o item a ).
Isso pode demostrar que a solução para o item b ) foi mais direta e menos interpretativa do que a do item a ), uma vez que a maior parte dos estudantes errou o entendimento da região na qual o torque é igual a 0, no eixo apresentado.
Por fim, uma pequena parcela conseguiu responder corretamente a todos os itens solicitados na questão.
De uma forma geral, a dificuldade de interpretação do sistema proposto pareceu ser sintomática e deve ser considerada para o aperfeiçoamento dos cursos de graduação na área de Engenharia Mecânica.
Poucas avaliações apresentaram interpretação correta do sistema mecânico proposto, tanto na parte de análise de torques quanto na parte de análise de reações em diferentes planos de ação.
Os estudantes apresentaram dúvidas no entendimento das regiões em que os torques são nulos no item a ), evidenciando uma possível falha no aprendizado referente ao conteúdo que envolve a análise de esforços em diferentes planos cartesianos, um tema comum nas atividades profissionais de um engenheiro mecânico.
Em relação às respostas dos estudantes, as estatísticas das correções demonstraram um fraco desempenho na questão avaliada.
Muitas respostas tiveram assertivas parciais e a grande maioria dos acertos estava relacionada ao item b ).
Por vezes, os estudantes apresentaram resultados sem qualquer desenvolvimento ou apresentaram leituras equivocadas do enunciado.
A ausência ou dificuldade em atender ao que foi solicitado no item a ) podem evidenciar uma fragilidade na compreensão crítica e interpretativa do sistema mecânico apresentado.
Respostas fracas não abordaram a questão ou atenderam apenas parcialmente ao item b ), por terem apenas aplicado a somatória das forças e momentos em diferentes planos, mas com erros interpretativos.
Respostas medianas, em sua maioria, atenderam integralmente ao item b ) e não atenderam ao item a ), ou atenderam parcialmente aos dois, indicando alguma familiaridade com o tema, porém com pequenos erros interpretativos de esforços em eixos e reações em planos cartesianos.
As poucas respostas boas atenderam integralmente ao item a ) e parcialmente ao item b ).
Uma minoria dos casos demonstrou atender integralmente à questão.
Dessa forma, a questão foi considerada de média dificuldade com um breve tempo necessário para sua resolução.
Esse fato, juntamente com a elevada frequência de notas com valor zero, seja pela questão deixada em branco ou por nenhum item acertado, leva ao entendimento que a grande maioria dos estudantes desconhece os conceitos de torque e reações atuantes nos sistemas mecânicos.
Assim, a média das notas corrigidas para esta questão discursiva foi fortemente impactada pelo elevado número de notas zero. 5.3.2.3 CONSIDERAÇÕES FINAIS SOBRE A QUESTÃO DISCURSIVA DO COMPONENTE ESPECÍFICO A questão apresentada para área de Engenharia Mecânica foi pertinente, redigida de forma acessível, abordou conteúdos fundamentais para a formação do profissional e apresentou uma dificuldade média com relação aos conceitos de análise de tensões atuantes em sistemas mecânicos de transmissão.
A análise global da questão revelou um desempenho fraco dos estudantes avaliados.
Em relação à questão avaliada, verificou -se um nível baixo de conhecimento.
Notou -se também erros na apresentação de unidades e análises dimensionais do sistema, sendo esse um erro grave, principalmente por se tratar de futuros profissionais em fase final de formação.
A maioria dos estudantes não conseguiu apresentar uma análise crítica, tampouco responder o que foi solicitado pelos itens da questão, pois poucos atentaram às informações apresentadas no desenho esquemático e aos itens questionados.
Alguns poucos demonstraram o domínio sobre o assunto ou mesmo um comprometimento e interesse na resolução da questão.
Apesar disso, é possível ponderar também que o tempo e a quantidade de questões podem ter gerado significativa influência no momento de elaborar corretamente a resposta ou mesmo na decisão de não responder, deixando a resolução em branco.

% matched lemmas: 
\end{textsample}
